\documentclass[10pt,a4paper]{article}
\usepackage[utf8]{inputenc}
\usepackage[vietnamese]{babel}
\usepackage{amsmath,amssymb,amsfonts}
\usepackage{geometry}
\geometry{a4paper, left=10mm, right=10mm, top=10mm, bottom=10mm}
\usepackage{tikz}
\usetikzlibrary{shapes,arrows.meta,positioning,calc}
\usepackage[most]{tcolorbox}
\usepackage{array,booktabs,multicol}
\usepackage{fontawesome}
\usepackage{helvet}
\renewcommand{\familydefault}{\sfdefault}
\usepackage{parskip}

\definecolor{primary}{RGB}{112, 48, 160}
\definecolor{teal}{RGB}{0, 128, 128}
\definecolor{orange}{RGB}{230, 108, 44}
\definecolor{green}{RGB}{46, 139, 87}
\definecolor{rose}{RGB}{192, 0, 0}

\pagestyle{empty}

\begin{document}

% --- HEADER BANNER ---
\begin{tcolorbox}[
  enhanced,
  colback=violet!8!white,
  colframe=violet!80!black,
  arc=2.5mm, boxrule=1pt,
  left=3mm, right=3mm, top=2mm, bottom=2mm
]
  \begin{tabular*}{\textwidth}{@{\extracolsep{\fill}} lr}
    {\small\bfseries\color{violet!85!black}\faBook\ SỔ TAY KIẾN THỨC TOÁN 2} & 
    {\small\bfseries\color{teal!85!black}\faGraduationCap\ CHƯƠNG TRÌNH GDPT 2018}
  \end{tabular*}
  \vspace{-1mm}
  \begin{center}
    \rule{0.98\textwidth}{0.6pt}
  \end{center}
  \vspace{-2mm}
  \centering
  {\LARGE\bfseries\color{violet!90!black}\faBookmark\ BÀI 58: LUYỆN TẬP CHUNG ĐẠI LƯỢNG ĐO LƯỜNG VÀ TIỀN TỆ}\\[1mm]
  {\itshape\footnotesize\color{black!75} Tổng ôn tập độ dài ($cm, dm, m, km$), khối lượng ($kg$), dung tích ($l$) và tiền tệ}
\end{tcolorbox}

\vspace{-2mm}

\noindent
\begin{minipage}[t]{0.49\textwidth}\vspace{0pt}
\begin{tcolorbox}[
  enhanced,
  colback=violet!4!white,
  colframe=violet!75!black,
  arc=2mm, boxrule=0.8pt,
  left=2mm, right=2mm, top=1.8mm, bottom=1.8mm,
  title={\footnotesize\bfseries\color{white}\faKey\ 1. Bảng tổng hợp các đại lượng đã học}
]
{\footnotesize
\begin{itemize}\setlength{\itemsep}{1pt}\setlength{\leftskip}{-2.5mm}
  \item \textbf{Độ dài:}
  \begin{itemize}\setlength{\itemsep}{0.5pt}\setlength{\leftskip}{-2mm}
    \item $1\,dm = 10\,cm$; $1\,m = 10\,dm = 100\,cm$; $1\,km = 1\,000\,m$.
  \end{itemize}
  \item \textbf{Khối lượng:} Ki-lô-gam ($kg$).
  \item \textbf{Dung tích:} Lít ($l$).
  \item \textbf{Thời gian:} Ngày, giờ, phút ($1\text{ ngày} = 24\text{ giờ}$; $1\text{ giờ} = 60\text{ phút}$).
  \item \textbf{Tiền Việt Nam:} $100$ đồng, $200$ đồng, $500$ đồng, $1\,000$ đồng.
\end{itemize}
}
\end{tcolorbox}
\end{minipage}
\hfill
\begin{minipage}[t]{0.49\textwidth}\vspace{0pt}
\begin{tcolorbox}[
  enhanced,
  colback=teal!4!white,
  colframe=teal!75!black,
  arc=2mm, boxrule=0.8pt,
  left=2mm, right=2mm, top=1.8mm, bottom=1.8mm,
  title={\footnotesize\bfseries\color{white}\faSliders\ 2. Kĩ năng giải toán có kèm đại lượng}
]
{\footnotesize
\begin{itemize}\setlength{\itemsep}{1pt}\setlength{\leftskip}{-2.5mm}
  \item \textbf{Quy tắc vàng 1:} Luôn đưa về cùng đơn vị đo trước khi so sánh, cộng, trừ.
  \item \textbf{Quy tắc vàng 2:} Thực hiện phép tính trên các con số, sau đó viết thêm tên đơn vị vào kết quả:
  $$35\,kg + 18\,kg = (35 + 18)\,kg = \mathbf{53\,kg}$$
  $$12\,l \times 3 = (12 \times 3)\,l = \mathbf{36\,l}$$
  \item Khi giải bài toán có lời văn: viết tên đơn vị vào trong dấu ngoặc đơn ở phép tính và đáp số đầy đủ.
\end{itemize}
}
\end{tcolorbox}
\end{minipage}

\vspace{-1.5mm}
\begin{tcolorbox}[
  enhanced,
  colback=orange!4!white,
  colframe=orange!80!black,
  arc=2mm, boxrule=0.8pt,
  left=2mm, right=2mm, top=1.5mm, bottom=1.5mm,
  title={\footnotesize\bfseries\color{white}\faSitemap\ 3. Sơ đồ tư duy các đại lượng đo lường thực tế lớp 2}
]
\centering
\begin{tikzpicture}[x=1.55cm, y=0.8cm]
  \tikzset{arr/.style={-{Stealth[scale=0.8]}}, every node/.append style={align=center}}
  \node[rectangle, rounded corners=2pt, draw=violet!80!black, fill=violet!20, font=\footnotesize\bfseries] (c) at (0,0) {ĐẠI LƯỢNG\\LỚP 2};
  \node[rectangle, rounded corners=2pt, draw=teal!80!black, fill=teal!15, font=\scriptsize\bfseries] (d) at (2.2,0.9) {Độ dài\\$cm, dm, m, km$};
  \node[rectangle, rounded corners=2pt, draw=orange!80!black, fill=orange!15, font=\scriptsize\bfseries] (k) at (2.2,0.3) {Khối lượng: $kg$\\Dung tích: $l$};
  \node[rectangle, rounded corners=2pt, draw=green!70!black, fill=green!15, font=\scriptsize\bfseries] (t) at (2.2,-0.3) {Thời gian\\Giờ, ngày, tháng};
  \node[rectangle, rounded corners=2pt, draw=red!75!black, fill=red!15, font=\scriptsize\bfseries] (m) at (2.2,-0.9) {Tiền tệ\\Đồng ($100 \dots 1\,000$)};

  \draw[arr, thick] (c) |- (d);
  \draw[arr, thick] (c) |- (k);
  \draw[arr, thick] (c) |- (t);
  \draw[arr, thick] (c) |- (m);
\end{tikzpicture}
\end{tcolorbox}

\vspace{-1.5mm}
\noindent
\begin{minipage}[t]{0.49\textwidth}\vspace{0pt}
\begin{tcolorbox}[
  enhanced,
  colback=green!4!white,
  colframe=green!60!black,
  arc=2mm, boxrule=0.8pt,
  left=2mm, right=2mm, top=1.8mm, bottom=1.8mm,
  title={\footnotesize\bfseries\color{white}\faCheckCircle\ 4. Bài tập thực hành tổng hợp}
]
{\footnotesize
\textbf{1. Tính:}\\
- $45\,m + 28\,m = \mathbf{73\,m}$.\\
- $60\,kg - 25\,kg = \mathbf{35\,kg}$.\\
- $5\,l \times 6 = \mathbf{30\,l}$; \quad $20\,l : 2 = \mathbf{10\,l}$.\\
\textbf{2. Bài toán lời văn:}\\
Mẹ mua một gói bánh giá $600$ đồng và một phong kẹo giá $300$ đồng. Mẹ đưa cô bán hàng $1\,000$ đồng. Hỏi cô bán hàng phải trả lại mẹ bao nhiêu tiền?\\
\textit{Giải:} Tổng số tiền mua bánh và kẹo là:\\
$600 + 300 = 900$ (đồng).\\
Số tiền cô bán hàng trả lại là:\\
$1\,000 - 900 = \mathbf{100}$ (đồng).
}
\end{tcolorbox}
\end{minipage}
\hfill
\begin{minipage}[t]{0.49\textwidth}\vspace{0pt}
\begin{tcolorbox}[
  enhanced,
  colback=rose!4!white,
  colframe=rose!70!black,
  arc=2mm, boxrule=0.8pt,
  left=2mm, right=2mm, top=1.8mm, bottom=1.8mm,
  title={\footnotesize\bfseries\color{white}\faExclamationTriangle\ 5. Lỗi học sinh hay quên}
]
{\footnotesize
\begin{itemize}\setlength{\itemsep}{1.5pt}\setlength{\leftskip}{-2.5mm}
  \item \textbf{Quên ghi tên đơn vị:} Kết quả tính là $73$ nhưng quên không viết chữ $m$ $\to$ bị trừ điểm.
  \item \textbf{Cộng các đơn vị khác loại:} Ví dụ $5\,kg + 3\,l$ là phép tính hoàn toàn \textbf{vô nghĩa}!
  \item Khi giải bài toán có hai bước tính, cần xác định rõ bước nào tính trước, bước nào tính sau.
\end{itemize}
}
\end{tcolorbox}
\end{minipage}


\end{document}
